(1) 电偶极辐射的横向远场可写 \(E_\theta(r,\theta,t)=-\omega^2p_0\sin\theta\cos[\omega(t-r/c)]/(4\pi\varepsilon_0c^2r)\)；符号取决于电偶极矩二阶导的相位约定。相应磁场 \(B_\varphi=E_\theta/c\)，使 \(\mathbf E\times\mathbf B\) 沿径向外指，并满足远场 \(|E|=c|B|\)。
(2) 平均径向坡印廷能流 \(\overline S_r=\varepsilon_0c\overline{E_\theta^2}=\omega^4p_0^2\sin^2\theta/(32\pi^2\varepsilon_0c^3r^2)\)，其中余弦平方周期平均为 \(1/2\)。乘以面积元 \(r^2d\Omega\) 得 \(d\overline P/d\Omega=\omega^4p_0^2\sin^2\theta/(32\pi^2\varepsilon_0c^3)\)。
(3) 角积分 \(\int\sin^2\theta\,d\Omega=2\pi\int_0^\pi\sin^3\theta\,d\theta=8\pi/3\)。所以 \(\overline P=\omega^4p_0^2/(12\pi\varepsilon_0c^3)\)，与半径无关，符合远场能量守恒；固定偶极幅度时频率翻倍使功率增加为原来的十六倍。
(4) 沿偶极轴 \(\theta=0,\pi\) 时横向辐射为零，赤道 \(\theta=\pi/2\) 时角功率最大；当 \(p_0\to0\) 或 \(\omega\to0\)，平均辐射功率趋零。离源较近的 \(r^{-2}\)、\(r^{-3}\) 场有周期往返交换的反应能量，不能用其瞬时能量密度直接等同 \(r^{-2}\) 的净辐射通量。
